\subsection{CP-06: API REST \texttt{/v1/api} de extremo a extremo}
\label{sec:cp06}

\FichaCaso
  {CP-06}
  {David Estevan Rodriguez Jurado}
  {09/11/2026}
  {API REST \texttt{/v1/api} de extremo a extremo}
  {Joan Emmanuel Orduz Chia}
  {11/11/2026}
  {Ejercitar las rutas de la capa REST con un cliente que no instala el SDK de servidor y comprobar que conserva el monto exacto, discrimina la transacción de la redirección, exige la pasarela en la consulta y agrupa los bancos PSE por pasarela. El caso usa la colección de Postman en local, peticiones con \texttt{curl}, el prototipo C como cliente real y un humo de Render.}
  {Funcional (sistema, caja negra) y de contrato de la API.}
  {\textbullet\ SDK compilado y \path{simulator-api/} con dependencias instaladas.\newline
   \textbullet\ Newman disponible con \texttt{npx newman}, o Postman con la colección y el ambiente \texttt{local}.\newline
   \textbullet\ El archivo \emph{cp06-pago.json} creado en la raíz del repositorio con el contenido de la Tabla~\ref{tab:cp06-peticiones}.\newline
   \textbullet\ Para el prototipo C: \texttt{npm ci} y \texttt{npm run build} en \path{prototypes/checkout-frontend-rest/}; el puerto 4100 libre.\newline
   \textbullet\ Para el humo de Render: acceso a Internet, el servicio despierto y el token de la API para la petición~8.}
  {Local. API de Simulación en \texttt{http://localhost:3000}; prototipo C en \texttt{http://localtest.me:4100} con \texttt{npm run start:local}. Los pasos 1 a 6 usan credenciales ficticias de cliente; el paso 7 usa el perfil del \path{.env}, cuyas llaves solo viajan al simulador local. El humo de Render usa \url{https://kit-pagos-colombia.onrender.com} y no cobra. Node.js 20.}

\subsubsection{Objetivo}

Comprobar que un cliente sin SDK puede crear un pago, consultarlo y listar bancos PSE por REST con la misma semántica que el SDK: monto exacto como cadena, estado normalizado y nativo, resultado que distingue \texttt{TRANSACTION} de \texttt{REDIRECT\_REQUIRED}, pasarela obligatoria en la consulta y bancos agrupados por pasarela. El caso falla si la capa REST aplana esas diferencias, si redondea el monto o si una ruta acepta una petición incompleta.

\subsubsection{Procedimiento}

\begin{TablaPasos}{Pasos del caso CP-06}
  1 &
  Levantar el simulador sin el perfil de Wompi del \path{.env}. &
  \texttt{cd simulator-api}; luego \path{WOMPI_PUBLIC_KEY= WOMPI_PRIVATE_KEY= WOMPI_INTEGRITY_SECRET= npm run dev}. &
  El simulador escucha en el puerto 3000 y \texttt{GET /health} responde 200. Con esas tres variables vacías, verifica la firma de integridad con el secreto por omisión, el mismo del ambiente \texttt{local} de Postman. &
  Sin perfil; 200. & Pasó & Simulador en 4817 \\ \hline
  2 &
  Ejecutar la carpeta de la capa REST con Newman. &
  Desde la raíz, la petición~1 de la Tabla~\ref{tab:cp06-peticiones}. &
  Las 4 peticiones de la carpeta \texttt{01. Capa REST Kit Pagos (/v1/api)} responden 2xx, pasa la aserción de \texttt{GET /v1/api/gateways} y Newman termina con código 0. &
  4 de 4 y 5 aserciones. & Pasó & Newman \\ \hline
  3 &
  Crear un pago de tarjeta de Wompi por REST. &
  Petición~2, con \emph{cp06-pago.json}. &
  Respuesta 201 con \texttt{outcome: "TRANSACTION"}; \texttt{transaction.amount} es \texttt{"50000.00"}, una cadena y no un número; \texttt{currency} es \texttt{COP}; la referencia se conserva; \texttt{status} y \texttt{rawStatus} son \texttt{PENDING}, porque Wompi crea la tarjeta pendiente. &
  201; monto como cadena; \texttt{PENDING}. & Pasó & \texttt{curl} \\ \hline
  4 &
  Consultar el estado con el identificador y la pasarela. &
  Petición~3, con el \path{transaction.gatewayTransactionId} del paso 3. &
  Respuesta 200 con \texttt{status} y \texttt{rawStatus} en \texttt{APPROVED}; el mismo identificador y el mismo monto del paso 3. El desenlace llega en la consulta, no en la creación. &
  200; \texttt{APPROVED}. & Pasó & \texttt{curl} \\ \hline
  5 &
  Verificar que la consulta exige la pasarela. &
  Petición~3 sin \path{?gateway=wompi}, y luego con \path{?gateway=paypal}. &
  Las dos responden 400 con \texttt{code: "INVALID\_REQUEST"}; la API no deduce la pasarela del identificador. &
  400 y 400. & Pasó & \texttt{curl} \\ \hline
  6 &
  Listar los bancos PSE de Wompi. &
  Petición~4. &
  Respuesta 200 con \texttt{pseBanks} de un solo elemento, \texttt{gateway: "wompi"}, y una lista de bancos con \texttt{code} y \texttt{name}; el primero es \texttt{\{ code: "1", name: "Banco que aprueba" \}}. &
  Un grupo; banco 1. & Pasó & \texttt{curl} \\ \hline
  7 &
  Cobrar con el prototipo C contra la API local. &
  Detener el simulador y levantarlo con \texttt{npm run dev}, con el perfil del \path{.env}; luego \texttt{cd} \path{prototypes/checkout-frontend-rest} \texttt{\&\& npm run start:local}, abrir \url{http://localtest.me:4100}, elegir Wompi y PSE, y pagar. &
  La página cobra contra \texttt{/v1/api} sin un servidor de pagos propio. La consola del navegador, con el prefijo \texttt{[tienda C]}, muestra las llamadas a \texttt{/v1/api}, la redirección al banco y la consulta posterior. El pago termina aprobado. &
  PSE aprobado. & Pasó & Captura \\ \hline
  8 &
  Ejecutar el humo de Render. &
  Peticiones~5, 6 y 7, sin token, y petición~8, con el token de la API; si el servicio está dormido, esperar a que despierte. Desde el 11/11/2026, Render exige el token en \texttt{/v1/api}; \texttt{/health} y \texttt{/docs} siguen abiertas, porque son de lectura y no tocan pasarelas ni dinero. &
  \texttt{/health} responde 200 y \texttt{/docs/json} devuelve el documento OpenAPI con su campo \texttt{openapi} en 3.0, los dos sin token. \texttt{/v1/api/gateways} responde 401 sin token y, con el token, 200 con las cuatro pasarelas y su \texttt{cardCapture}. Si Render no está disponible, el paso se registra como bloqueado y no invalida la corrida local. &
  200 en los tres, sin token. & Pasó & Render. Corrida anterior a las 20:41, cuando Render aún no exigía el token \\ \hline
\end{xltabular}
\end{TablaPasos}

Cada petición es una sola línea de shell. El PDF la parte por el ancho de la página, así que para pegarla en la terminal se copia de \path{docs/final-doc/pruebas/sections/casos-prueba/cp-06.tex}, donde está en una línea.

\begin{table}[H]
\centering
\footnotesize
\renewcommand{\arraystretch}{1.15}
\begin{tabularx}{\textwidth}{|>{\raggedright\arraybackslash}p{2.6cm}|>{\raggedright\arraybackslash}X|}
\hline
\textbf{Petición o archivo} & \textbf{Contenido} \\ \hline
\emph{cp06-pago.json} & \path{{"gateway":"WOMPI","amount":50000,"currency":"COP","description":"Pago de prueba CP-06","orderReference":"ORD-API-CP06-1","payer":{"email":"postman.tester@example.com","fullName":"Tester Postman","documentType":"CC","documentNumber":"123456789","phoneNumber":"3001234567"},"paymentMethod":{"type":"CARD","cardToken":"tok_test_card_postman"}}} \\ \hline
Petición~1 & \path{npx newman run simulator-api/postman/kit-pagos-simulator.postman_collection.json -e simulator-api/postman/local.postman_environment.json --folder "01. Capa REST Kit Pagos (/v1/api)"} \\ \hline
Petición~2 & \path{curl -s -X POST http://localhost:3000/v1/api/payments -H 'Content-Type: application/json' -H 'x-gateway-public-key: pub_test_wompi_key' -H 'x-gateway-private-key: prv_test_wompi_key' -H 'x-gateway-integrity-secret: test_integrity_kit_pagos_simulator' -d @cp06-pago.json} \\ \hline
Petición~3 & \path{curl -s 'http://localhost:3000/v1/api/payments/<id>?gateway=wompi' -H 'x-gateway-public-key: pub_test_wompi_key' -H 'x-gateway-private-key: prv_test_wompi_key'} \\ \hline
Petición~4 & \path{curl -s 'http://localhost:3000/v1/api/pse-banks?gateway=wompi' -H 'x-gateway-public-key: pub_test_wompi_key' -H 'x-gateway-private-key: prv_test_wompi_key'} \\ \hline
Petición~5 & \path{curl -s https://kit-pagos-colombia.onrender.com/health} \\ \hline
Petición~6 & \path{curl -s https://kit-pagos-colombia.onrender.com/v1/api/gateways} \\ \hline
Petición~7 & \path{curl -s https://kit-pagos-colombia.onrender.com/docs/json} \\ \hline
Petición~8 & \path{curl -s https://kit-pagos-colombia.onrender.com/v1/api/gateways -H 'Authorization: Bearer <token de la API>'} \\ \hline
\end{tabularx}
\caption{Archivo y peticiones del caso CP-06}
\label{tab:cp06-peticiones}
\end{table}

\subsubsection{Criterio de aceptación}

\begin{enumerate}
  \item Las cuatro peticiones de la carpeta \texttt{01} pasan en el ambiente \texttt{local}.
  \item El monto se conserva como cadena exacta \texttt{"50000.00"}; la API no lo convierte a número ni lo redondea.
  \item La creación distingue \texttt{TRANSACTION} de \texttt{REDIRECT\_REQUIRED} con el campo \texttt{outcome}, y la consulta devuelve el estado normalizado y el nativo.
  \item La consulta sin \texttt{gateway} y con una pasarela desconocida responde 400 \texttt{INVALID\_REQUEST}.
  \item \path{/v1/api/pse-banks?gateway=wompi} devuelve un solo grupo, con bancos que tienen \texttt{code} y \texttt{name}.
  \item El prototipo C completa un pago PSE contra \texttt{/v1/api} sin un servidor de pagos propio.
  \item El humo de Render responde, o su indisponibilidad queda registrada como bloqueo ajeno al caso local.
\end{enumerate}

\subsubsection{Observaciones de diseño}

El caso elige la tarjeta de Wompi y no un PSE para la comprobación principal porque la creación con tarjeta devuelve \texttt{TRANSACTION} y obliga a consultar para ver el desenlace; así se ejercitan las dos rutas de lectura en una sola corrida. El PSE aparece en el prototipo C, donde el cliente sí necesita ver la redirección y la diferencia entre \texttt{TRANSACTION} y \texttt{REDIRECT\_REQUIRED}.

El paso 1 vacía las tres variables de Wompi porque, con un perfil completo en el \path{.env}, el simulador verifica la firma de integridad con el secreto de ese perfil y no con el valor por omisión que trae el ambiente \texttt{local} de Postman (\path{simulator-api/postman/README.md}). Así el caso no necesita copiar un secreto real en ningún comando. El paso 7 vuelve al perfil del \path{.env}, porque C no envía credenciales y la API usa las del servidor.

El humo de Render se mantiene separado de la aceptación local. Render es un servicio del nivel gratuito que puede tardar en despertar, y una caída del despliegue no dice nada sobre la corrección funcional de \texttt{/v1/api}; por eso se marca como bloqueado, no como fallo del caso. Render no se usa para la carga.

\subsubsection{Trazabilidad}

RF-13 (OpenAPI), RF-16 (capa REST), RF-18 (motivo de rechazo serializado) y RF-19 (bancos PSE); HU-09, HU-11 y HU-12.
